\section{Limitations}
\label{sec:limitations}

\textbf{The MMLU deficit is real, reproduced, and unresolved.} $-14.6$
points against FP16 on the 4B served arm, reproduced three times on two
machines; $-14.73$ paired on the sealed artifact, the arm all paired
statistics use; 10.57 at 8B, 6.85 at 14B, all three resolved. Abstract
algebra and accounting fall to 25\%, chance, where history, law and
psychology hold above 80\%: two bits damage \emph{reasoning} far
more than \emph{recall}, mostly what a perplexity corpus measures, so
perplexity alone is never evidence of quality at this rate. Suspects,
measured and bounded: output-side rotation (near-zero marginal effect in the
original paper's ablation); a free-magnitude variant with a closed-form
scale solve ($\times1.99$ perplexity at full depth on the 0.6B proxy, killed
by a pre-registered gate); and calibration volume --- our ${\sim}131$k
tokens are roughly $100\times$ less than the original paper's, but a
deliberate-contamination oracle (calibrating on the evaluation text itself,
an upper bound on what volume or corpus buys) recovers only $1.6\%$ of
perplexity on a proxy, a $\times13$ volume curve $1.2\%$; a bound weaker
than it looks (perplexity, 3 blocks, a 0.6B proxy; it does not bound MMLU),
but enough that we did not commission a $\times100$ run. Untested:
reasoning-weighted calibration composition, learned per-column scales (about
two MMLU points under the original paper's fine-tuning protocol), low-rank
compensation adapters.

\textbf{We are slower than the 2-bit kernel we set out to beat, and the
margin is large.} QTIP's kernel runs $2.27\times$ our production layout on
matching shapes in the same process (\S\ref{sec:qtip}). We measured it as
shipped, in its own launch geometry, tuning nothing either way: a tuning
pass could move that number, and we did not attempt one. Not a defect of
our implementation: both formats store the same rate on
disk, and the gap is entirely the load-time unfolding a combinatorial
lattice index requires and a sliding-window trellis does not. A
research direction, not a patch, and not obviously open: the granular gain
$\Lambda_{24}$ contributes over a deep trellis is at most a decibel, which a
long-memory trellis on incoherence-processed (hence approximately i.i.d.\
Gaussian) weights already captures --- the reading our measurements support,
not a proof. The same gap remains for Marlin's batched 4-bit GEMM and the
other 2-bit vector-quantized systems of \S\ref{sec:related}.

\textbf{Three models end to end, one serving GPU, greedy decoding.}
End-to-end: Qwen3-4B, 8B and 14B on one L40S at batch size 1;
only the layout bench runs on a second architecture, where it bounds rather
than extends the claim (\S\ref{sec:a100}). Batching changes
the arithmetic-intensity regime that makes the kernel memory-bound; larger
models change the embedding share and tail-column overhead.

\textbf{The scale argument is three points, and not a law.} On the 4B,
4-bit AWQ dominates every axis but disk and served footprint
(Table~\ref{tab:campaign}), a ${\sim}3\%$ b/param margin that does not,
alone, justify a 14.6-point MMLU loss. The 8B moves every relevant
number in the 2-bit direction (Table~\ref{tab:campaign8b}); the 14B keeps it
and \emph{loses the rate} (Table~\ref{tab:scale}). Whether it
continues to the 70B class --- the regime the sovereignty argument needs ---
is measurable with this methodology, and not measured here; the third
point is why we do not extrapolate. Those qualifications
(one verdict per metric on the slowdown; intervals paired, on committed raw
output; a memory margin at all three sizes, not monotone) are carried once,
in \S\ref{sec:evaluation} and Appendix~\ref{app:scale}.

\textbf{Our speedups are over FP16, not over the state of the practice at
either rate.} At the \emph{kernel} level both comparisons go against us:
the AWQ GEMV converts 88\% of its byte bound against our 65\%, and QTIP is
$2.27\times$ faster outright (\S\ref{sec:layouts}, \S\ref{sec:qtip}). At
the \emph{engine} level the comparison exists only within each stack
($2.413\times$ for AWQ--Marlin in vLLM against $1.11\times$ for us in
\texttt{candle}; \S\ref{sec:integration}), and goes against us too --- with
the caveat that the stacks differ independently of weight format:
vLLM's FP16 control reaches 83.1~tok/s where ours reaches 43.6, so roughly
a factor 1.9 of the end-to-end distance is engine, not decoder. A cost to
the user and a confound to the reader: we report the two ratios side by
side and never divide them. Tables~\ref{tab:campaign}
and~\ref{tab:campaign8b} carry no 4-bit throughput column: the two tok/s are
different measurements.

\textbf{Speedups carry ranges, and the two protocols differ.} Inter-process
dispersion on identical binaries spans several percent, so every bench
ratio (\S\ref{sec:layouts}, Table~\ref{tab:layouts},
Figure~\ref{fig:layouts}) is a median of per-round ratios over interleaved
arms, with its min--max range. \S\ref{sec:integration}'s end-to-end
figures are medians of five timed generations per arm, but
their ratios are quotients of two medians with a conservative envelope: the
arms load their models exclusively and their rounds never coexist.
Effects smaller than a range are unresolved; second decimals read as
dispersion, not measurement.

\textbf{The rotated basis is a real serving constraint.} Weights sit in the
incoherence-rotated basis, so any consumer must apply the activation
rotation. Our rotation kernel is verified bit-exact on CUDA and
exists nowhere else (no Metal port yet), so the served path is currently
NVIDIA-only ---
awkward for the sovereignty motivation, and first in line for future work.

\textbf{That same rotation caps the servable width one model class below
the one the argument wants.} A Walsh--Hadamard transform's barrier
ladder, logarithmic in the transform length, synchronizes one block, so the
whole activation must sit in that block's shared memory: one f32 per
coordinate whatever the input dtype. The binding width is
\texttt{intermediate\_size}, the input of \texttt{down\_proj}; the ceiling
is arithmetic:

\begin{center}\footnotesize
\setlength{\tabcolsep}{4pt}%
\begin{tabular}{lrrll}
\toprule
Model & \texttt{intermediate\_size} & shared bytes & vs.\ 49\,152 default & vs.\ 101\,376 opt-in \\
\midrule
Qwen3-4B  &  9\,728 &  38\,912 & fits & fits \\
Qwen3-8B  & 12\,288 &  49\,152 & fits \emph{exactly} & fits \\
Qwen3-14B & 17\,408 &  69\,632 & does not fit & fits \\
Qwen3-32B & 25\,600 & 102\,400 & does not fit & \textbf{short by 1\,024 bytes} \\
\bottomrule
\end{tabular}
\end{center}

\noindent
This hardware exposes \emph{two} bounds and our host code checked one: a
guard against the 49\,152-byte per-block default refused to launch the 14B
--- correctly, an overrun would corrupt silently --- costing one job
(\$0.24, no tokens produced) and finding the wall. The larger bound is
opt-in, requested once per function, and \emph{measured} on the card that
served (101\,376 bytes); only the host arithmetic changed.
The 8B sat on the default bound to the byte, unflagged because nothing ever
failed. The 32B misses the opt-in bound by 1\,024 bytes, the driver's
reserve, which no host-side trimming reclaims. The bound holds for any
format: every layout of \S\ref{sec:layouts} calls the same rotation. One
escape is named and
deliberately \emph{not} designed: half-precision shared activations would
halve the footprint below the opt-in bound, but at $n = 25\,600$ that puts
an accumulation of depth equivalent to ${\sim}14.6$ Hadamard stages in half
precision --- a trade we
have neither made nor priced.
